\section{Methods}\label{sec:methods}
\subsection{Bot Implementation}\label{sec:bot-implementation}

Eddie Bot is implemented as a Flask service backed by Azure App Service. It listens to mentions of a single X account via the Filtered Stream API from X, accepts invocations from enrolled participants only, and processes each accepted mention through a seven-stage verification pipeline before posting a reply. The bot is operated under an IRB-exempt determination, and the full source code is released alongside this report. In this section, we describe some of the key design choices in the pipeline.

\paragraph{Multimodal extraction.} When the parent post contains an image, the bot invokes a vision-language model to extract OCR text, a structured description, and a \textit{canonical-image match} flag indicating whether the image is a known artifact (e.g., iconic news photographs that frequently recirculate decontextualized). Videos and animated GIFs are handled via their X-API \texttt{preview\_image\_url} thumbnail, which the VLM treats identically to a still.

\paragraph{Claim extraction and action selection.} A single LLM call decomposes the parent post into atomic propositions, marks exactly one as the \textit{central} claim, and selects one of five \textit{actions} for the mention: \texttt{verify} (the central claim is falsifiable), \texttt{provide\_context} (the literal claim is true but the framing omits something material), \texttt{challenge\_opinion} (the claim is a strongly-stated opinion on which credible critics have published push-back), \texttt{surface\_perspectives} (the topic is contested), or \texttt{decline} (no actionable angle). The selection conditions on both the claim's character and any instruction the invoker wrote in the mention. When the invoker's stated ask does not fit the claim (e.g., a request to ``fact-check'' a pure opinion), the pipeline pivots to the action that does fit and notes the pivot in the final message.

\paragraph{Iterative verification.} For non-decline actions, the bot queries the public web iteratively, conditioning each search question on the results of prior queries. Retrieved pages are passed through a clean-extraction layer that returns the main article body (stripping navigation, ads, and boilerplate), so downstream reasoning operates on the actual reporting rather than search-result snippets. Evidence is deduplicated by URL across questions to avoid double-counting a single source.

\paragraph{Reconciliation under a source-quality check.} Retrieved evidence is partitioned by stance (supports, refutes, disputes, or contextualizes the central claim), and each URL is classified into one of six tiers (\textit{fact-checker}, \textit{reputable news}, \textit{primary source}, \textit{aggregator}, \textit{low-quality}, or \textit{satirical}). Classification uses a curated table of roughly 910 domains assembled from three sources: the IFCN's published list of verified signatories (mapped to the \textit{fact-checker} tier), Wikipedia's perennial-sources list (whose per-source reliability ratings supply the \textit{reputable news} and \textit{low-quality} tiers), and a hand-curated editorial supplement for categories no public list covers (\textit{primary source} and \textit{satirical} domains, together with established news outlets and fact-checkers not currently present in the two public lists). Each list is pinned to a specific upstream version (the IFCN commit and the Wikipedia revision identifiers) and recorded in every frozen verdict for reproducibility, and domains absent from the table fall to a model-based classifier. Clearly non-evidentiary domains such as stock-image libraries, shopping sites, and user-generated-content farms (61 in total) are dropped before classification so they never enter the evidence base. A confidence gate then determines whether the bot may commit to a finding: the \texttt{verify} and \texttt{provide\_context} actions require at least two \textit{distinct} reliable-tier sources before the verdict can be reported; \texttt{challenge\_opinion} and \texttt{surface\_perspectives} require at least one reliable-tier citation per counterpoint or perspective. When the gate fails, the action outcome is downgraded to a structured \textit{no usable evidence} form, and the renderer reports that state rather than the original claim.

\paragraph{Audit and Freeze.} A post-hoc audit verifies that the reconciled payload is internally consistent: every URL appearing in the text or in any structured citation field must also appear in the source-quality table, every action carries the structural fields its outcome requires (e.g., a refuted verdict requires a counter-fact), and the central claim must occupy exactly one stance bucket. On audit failure, the verdict is forced to \textit{NotEnoughEvidence} and any reconcile-stamped evidence stances are reset to neutral, so the frozen record cannot carry contradictory signals between its findings and its supporting evidence. The full verdict (atomic claims, retrieved evidence with stances and clean article bodies, the per-mention source-quality table, the chosen action and its outcome, the reconciled findings, and the presentation payload that drives rendering) is serialized to disk as an immutable research artifact. Each frozen verdict is the unit of analysis for the audit and methods-reproducibility claims in this paper.

\paragraph{Render Reply.} The reply tweet is composed by combining a per-action template (one of the five actions above) with a per-tone register (\textit{agreeable}, \textit{neutral}, or \textit{satirical}) and a fixed set of hard constraints: no URLs in the reply body, sources named only by their display name (e.g., ``Snopes'', ``AP News''), and an explicit yes/no opener when the invoker posed a yes/no question. The dossier link is then posted as a separate self-reply pointing to a per-mention \texttt{/info} page that lists the retrieved sources and the reconciliation reasoning. This split keeps the reply itself compact and free of URLs while keeping the full evidence trail one click away. For the audit, it also makes each posted reply individually traceable to its underlying evidence. The full set of render-stage prompts (the five action templates, three tone registers, and the shared hard constraints) is reproduced in \autoref{appendix:render-prompts}.

\subsection{Field Study}\label{sec:field-study}
To compare the effects of our different experimental conditions, we deploy Eddie Bot in a between-subjects randomized controlled trial on X (previously Twitter). We recruit paid participants to serve as ``invokers,'' who are instructed to mention Eddie Bot on 10 posts a day over a 5-day period. To mitigate the risk of invalid participant behaviors (such as invoking the bot on random posts), we recruit participants through trusted channels at our university and social networks. We recruit 30 participants in total.
% \todo{Add summary table of participant demographics later}

At the beginning of the study, invokers are assigned to a bot condition (agreeable, neutral, satirical) and a political topic (LGBT issues, immigration, healthcare, economy, religion, and race). To search for controversial posts, we provide invokers a search link each day of the study that filters for posts matching topic-related keywords, written in English, and with at least 20 replies. Invokers are instructed to select at least 10 posts from the returned posts, or the replies to the returned posts, that they disagree with and would like to fact-check. They invoke Eddie Bot by mentioning Eddie Bot in a comment and including a short contextualizing instruction (e.g. ``\texttt{@eddiexbot Is this true?}''). Once invoked, Eddie Bot replies to the invoker with a publicly visible reply generated via the pipeline described above, as well as a link to an informational page listing the sources and reasoning behind the reply. Other X users who are not recruited as part of the study (``bystanders'') can then engage with Eddie Bot's posts. This study design was approved by our institution's IRB and invokers were compensated \$60.

We collect several metrics of interest to capture effects on \textit{engagement} and \textit{affective polarization} across both invokers and bystanders. Before and after participating in the study, invokers are asked to fill out a survey measuring their agreement with statements corresponding to different types of affective polarization. Both surveys also ask questions about invokers' comfort with engaging in fact-checking behavior on X, and the post-survey asks invokers questions about their overall experience using the bot. In addition, invokers are sent a daily survey during the study period asking them to rate the quality of the bot's replies from the previous day. We interpret these responses as a proxy for engagement, assuming that a higher quality rating corresponds to a greater willingness to continue using the bot. We capture bystander engagement by recording the number of likes, shares, replies, link clicks, and views on Eddie Bot's posts. Because Eddie Bot posts its own dossier link as a self-reply to each fact-check reply (\autoref{sec:bot-implementation}), and X counts that self-reply in the parent's public reply count, we exclude the bot's self-reply from both the reply count and the corpus of collected replies, so that all reply-based measures reflect only genuine bystander replies. We capture affective polarization by applying two LLM-based classifiers for sentiment and signs of antagonistic affective polarization on these bystander replies to Eddie Bot's posts. See \autoref{appendix:experimental-details} for more details.